\documentclass[12pt]{article}
\usepackage[utf8]{inputenc}
\usepackage[T1]{fontenc}
\usepackage{hyperref}
\title{Yevanic Language Resources for K-12 Curriculum Development}
\author{Compiled by Claude Sonnet 5.5 (claude-sonnet-5-5) for Andr\'as Kornai}
\date{}
\begin{document}
\maketitle

\section{Introduction}
Yevanic (ISO 639-3 \texttt{yej}, Glottocode yeva1238), also called Judaeo-Greek, Romaniote or Yevanitika, is the Jewish Greek variety of the Romaniotes and the Constantinopolitan Karaites. It is written in the Hebrew alphabet. According to Wikipedia, only a few semi-speakers were known in Israel in 1987, and there may be a handful of elderly speakers in Turkey. The Endangered Languages Project rates it severely endangered.

\section{Primary Educational Resources}
No primers, courses or dictionaries were found in the searches.

\section{Linguistic and Reference Materials}
\begin{itemize}
\item Wikipedia article on Yevanic: \url{https://en.wikipedia.org/wiki/Yevanic}
\item Endangered Languages Project resources page for Yevanic. It lists a bibliography, media, and Krivoruchko's ``Judeo-Greek'' entry from the Jewish Language Research Website: \url{https://endangeredlanguages.com/node/38751/resources}
\item Endangered Languages Project language page: \url{https://www.endangeredlanguages.com/node/38751}
\end{itemize}

\section{Religious and Community Resources}
The searches did not find any dedicated community or religious resources.

\section{Digital and Audio Resources}
The ELP page links manuscript examples, including an excerpt from Abraham Ibn Ezra's astrological treatise at the Bodleian Library and a tenth-century Hebrew-Greek glossary in Leningrad. These are historical texts, not teaching audio.

\section{Curriculum Development Context}
Little exists for teaching Yevanic. A curriculum would have to rely on the ELP bibliography and the historical manuscripts. A heritage-oriented unit linked to Romaniote Jewish history would be the realistic starting point.
\end{document}